\backchapter{Epilogue}

\noindent The book closes with an epilogue in two voices. The first, new in this edition, is by the
reviser; the second is Shelley's own closing survey of the whole story.

\ssection{Epilogue, by R. L. Hatchett}
Hatchett begins with a quotation from Gordon Fee's \emph{Paul, the Spirit, and the People of God}. To
be saved, in Paul's view, means to become part of the people of God, born by the Spirit into God's
family, joined to one another as one body and formed into God's temple. God is not simply saving
individuals and preparing them for heaven. He is creating a people for his name, among whom he can
dwell and who in their life together reproduce his life and character in its unity and diversity.

The book, Hatchett says, has left a trail of stories that explain who ``we'' are: the right-believing
early Christians who kept the church on course by grasping something of the mystery of who Jesus is;
the Protestant Reformers who sought the heart of the gospel in Scripture; the believers who held to
the faith instead of yielding to liberal leanings; and the awakenings that shaped a prosperous
culture. The profile that results (orthodox, Protestant, conservative or evangelical, North American)
fits many of Shelley's students and readers over the years.

It is right to tell our own story and find our place in God's grand story. But a preoccupation with
its success may be costly. A story admired too much begins to look like the summit or the conclusion.
Then we fail to see and value others and to measure ourselves rightly, and with our story wrongly
appraised we grow discouraged and lose track of where God's story is heading. The unprecedented
growth of Christianity in the last hundred years calls for reassessment. The two ends of Christian
history now speak in unison. The first Christians, as Fee recalls, understood that God was forming a
people drawn from all the peoples of the world, often from low rank, and given the privilege of
belonging to him. Today God is again gathering a people from almost every conceivable place. The
hardships of the Western church may be read partly in this light, as a microcosm of what is happening
in the world: Western Christians are losing status and even facing persecution, and at the same time
they are witnessing the ingathering of this diverse people.

Problems and dangers abound, and it is not irreverent to ask questions. Are the new believers
superstitious or naive? Is their faith nominal, paper thin? Can they break with their culture? How
could the genocide in Rwanda happen on the ground of the East African Revival? But the same kind of
question must be asked of the West. Is our practice of Christianity shaped or misshaped by an
affluent culture? Did not Christian brother bear arms against Christian brother in the American Civil
War? And there are things to learn from the Global South. Do we understand solidarity? Is our
metaphysics stifling to the Spirit? Appraising our stories again may lead us to value belonging to
the people of God above everything else.

A circle persists: reading church history rightly helps us read the Bible rightly, and reading the
Bible rightly helps us read church history rightly. Both gain when we discern the theological
significance of the church as God's people gathered from the nations. Hatchett ends by quoting 1
Peter 2:4--10 in full. Believers come to Christ, the living stone rejected by mortals but chosen and
precious in God's sight, and are built like living stones into a spiritual house and a holy
priesthood. They are ``a chosen race, a royal priesthood, a holy nation, God's own people,'' called
out of darkness into his marvelous light: ``Once you were not a people, but now you are God's
people.''

\ssection{Epilogue, by Bruce Shelley}
After two thousand years Christianity is, at least nominally, the faith of one-third of the earth's
population. From a handful of fishermen, tax collectors and youthful troublemakers in an obscure
province it has spread over the globe to claim the loyalty of nearly a thousand million people. Yet
few of these professed believers have seriously studied the history of their religion. In earlier
ages the adherents of one faith seldom met those of another, and few had to defend their religion
against a rival. Now that mass media make the world a neighborhood, such ignorance is hard to
justify. The separation of church and state has all but removed religion from public education, and
Christian education in many denominations has given members no adult understanding of the faith. It
is no surprise, then, that today's Christians so often blend gross errors with orthodox confession or
defend some pagan practice as Christian conduct.

Informed Christians may be tempted to despair, but they know that human failure is only half the
story. The church has often been its own worst enemy, and renewal has often come from a wholly
unexpected source. Again and again some unseen Power has turned aside a threat or made a crisis into
an opportunity. Persecution purged the household of faith, the spread of heresy clarified the
church's basic beliefs, and the sudden arrival of the barbarians opened doors for expansion. This
capacity to face new challenges and tap the sources of renewal is one of the secrets of
Christianity's growth. The way forward has usually meant a careful look backward, to the image of
God revealed in the story of Jesus.

Shelley then retells the whole story in brief, age by age. Christians have always treated the age of
Jesus and the apostles as a model for all the others. It gave the church its faith in Jesus, the
risen Messiah, and the hope of forgiveness through him, and it showed in the life of Paul that the
gospel of grace knows no boundary of nation, race, sex or culture. The catholic Christianity that
accepted this truth spread quickly through the Mediterranean world. It confronted Gnosticism,
Marcionism and Montanism and called a lie a lie by appealing to the apostolic writings and to the
orthodox bishops who guarded them, while its martyrs faced the persecuting power of Rome and died as
witnesses for others to follow.

The seed of the martyrs' blood, as Tertullian called it, bore fruit in the conversion of the empire.
The Imperial Age began in 312, when Constantine caught a vision of Christ, and before the fourth
century ended Christianity was the official religion of the Roman state. A church in the catacombs
was one thing, but what has Christianity to do with palaces? Under the emperors' tutelage the church
learned to serve the seats of power by formulating the faith for the masses, in the age of the great
councils. Those with no taste for palaces went to the wilderness to seek another way to grace, and
the revered hermits found themselves at the head of monasticism, the wave of the future. Most
Christians saw the hand of God in the marriage of church and state. In the East it lasted a thousand
years, and a mystical piety flourished under orthodox emperors until 1453, when the Muslim Turks
brought the Byzantine Empire to its final ruin. The fall of Constantinople meant the rise of Moscow,
the new capital of Eastern Orthodoxy.

In the West, after the Germans and Huns broke the empire's defenses in the fifth century and entered
the eternal city itself, people turned to Augustine's \emph{City of God} and found in it a vision for
a new age. We call those centuries medieval; the people who lived in them called them Christian. The
reason lies in the pope, who stepped into the ruins of the Western empire and built the medieval
church on Rome's past glory. As the one surviving link with that past, the Church of Rome sent out
Benedictine monks as missionary ambassadors to the German peoples. Over centuries the popes, aided
by Christian princes, pacified and baptized a continent and called it Christendom. But baptized
masses meant baptized pagans, and by the tenth century renewal was plainly necessary. It began at
the monastery of Cluny in central France and spread until it reached the papacy. The greatest of the
reforming popes was Gregory VII, and his zealous successors carried the office to the height of
earthly power. The church of the twelfth century was itself a kind of empire, a spiritual and
earthly kingdom reaching from Ireland to Palestine and from earth to heaven, as the crusades and
Scholastic philosophy bear witness. Power corrupts, however. The church gained the world and lost its
soul, or so a steady stream of reformers preached: Waldenses, Franciscans, Albigenses. Amid the
strife for power and the barren religion of the fourteenth and fifteenth centuries many turned to
the Bible for fresh vision.

Reform came with a fury. Luther sounded the trumpet and hosts rallied to it. The Reformation was the
mobilization of Protestantism: Lutheran, Reformed, Anglican and Anabaptist. By the mid-sixteenth
century it had shattered the unity of western Europe and bequeathed religious pluralism to modern
times. Rome resisted. It mustered new troops, above all the Society of Jesus, sent fresh waves of
missionaries to Asia, Africa and Latin America, and waged war in France, the Netherlands and Germany.
In the end Christendom slipped into the past, and in its place came the denominational idea of the
church, which let modern nations treat the churches as voluntary societies separate from the state.

New schools of thought filled the seventeenth century, none stronger than Reason, which asked who
needs God, since man can manage alone. Christians protested loudly, but the idea spread until
secularism filled the public life of Western societies and God remained only as a matter of personal
choice. Unable any longer to call on the arm of power against such heresies, many Christians turned
to the way of the apostles, prayer and preaching. The result was the evangelical revivals, chiefly
Pietism, Methodism and the Great Awakening, which tried through preaching and personal conversion to
restore God to public life.

In the Age of Progress Christians of all kinds fought the advance of secularism. From the evangelical
awakenings came new efforts to carry the gospel to distant lands and a host of social ministries in
industrial Europe and North America, while from the ramparts of Rome a defensive papacy fired at the
modern enemies of the Catholic faith. Despite their best efforts Christianity was slowly driven from
public life in the West, and believers were left with the problem of our own time: how can Christians
exert moral influence in pluralistic and totalitarian societies where Christian assumptions about
reality no longer prevail? The depth of the problem appeared in the Age of Ideologies, when new gods
claimed the loyalty of secular people. Nazism exalted the state, Communism worshiped the party, and
American democracy revered the individual and his rights. Supposedly enlightened nations fought two
global wars for the supremacy of these deities, and when none prevailed a cold war of coexistence
settled on the once Christian nations. Through these years the denominations struggled over orthodox
and liberal theologies, sought ways to recover a lost unity and showed a new hunger for apostolic
experience.

After the Second World War vigorous new Christian leadership emerged in the Third World and offered
hope of a new day for the old faith. Had missionaries from the neo-pagan nations of Europe and North
America given Christianity a stake in the future by carrying the gospel to Asia, Africa and Latin
America? Only time will tell. But Christians can hope, because faith reaches beyond earthly
circumstances and rests in a person. No other person in recorded history has influenced so many
people in so many conditions for so long as Jesus Christ. The shades of his image shift with human
need: the Jewish Messiah of the believing remnant, the Wisdom of the Greek apologist, the Cosmic King
of the imperial church, the Heavenly Logos of the orthodox councils, the World Ruler of the papal
courts, the monastic Model of apostolic poverty, the personal Savior of the evangelical revivalists.
He is, Shelley concludes, a man for all time. In a day when many think him irrelevant, a relic of a
discarded past, church history gives quiet testimony that Jesus Christ will not disappear from the
scene. ``His title may change, but his truth endures for all generations.''
